% manuel_recouvrement_suite.tex — inclus par manuel_recouvrement.tex

% -------- Onglet Cat\'egories --------
\section{Onglet 1 --- Cat\'egories de Variables}

Les cat\'egories permettent de regrouper les frais par famille logique. Exemples~: \textit{Minerval}, \textit{Fournitures}, \textit{Activit\'es parascolaires}.

\begin{etape}[title=Proc\'edure : Cr\'eer une cat\'egorie]
\begin{enumerate}[leftmargin=1.5cm,label=\textbf{\arabic*.},itemsep=10pt]
  \item \textbf{Cliquez sur l'onglet} \onglet{Cat\'egories} dans la barre de configuration.
  \item \textbf{Remplissez le champ} \champ{Nom} dans le formulaire en haut (ex.\ : \textit{Minerval}).
  \item \textbf{Cliquez sur} \bouton{Ajouter}. La cat\'egorie appara\^it dans le tableau ci-dessous.
\end{enumerate}
\end{etape}

\begin{resultat}
La cat\'egorie est imm\'ediatement disponible dans le menu d\'eroulant \champ{Cat\'egorie} lors de la cr\'eation de variables (onglet suivant).
\end{resultat}

% -------- Onglet Variables --------
\section{Onglet 2 --- Variables de Paiement}

Les variables repr\'esentent les \'el\'ements concrets \`a payer. Chaque variable est rattach\'ee \`a une cat\'egorie.

\begin{etape}[title=Proc\'edure : Cr\'eer une variable]
\begin{enumerate}[leftmargin=1.5cm,label=\textbf{\arabic*.},itemsep=10pt]
  \item \textbf{Cliquez sur l'onglet} \onglet{Variables}.
  \item \textbf{S\'electionnez la cat\'egorie} dans le menu \champ{Cat\'egorie}.
  \item \textbf{Saisissez le nom} dans \champ{Nom variable} (ex.\ : \textit{Frais scolaire T1}).
  \item \textbf{Cliquez sur} \bouton{Ajouter}. La variable appara\^it dans le tableau.
\end{enumerate}
\end{etape}

\begin{info}
Chaque variable dispose d'un bouton \textbf{Obligatoire} (Oui/Non) qui permet de marquer les frais obligatoires. Cliquez sur le bouton vert/orange dans la colonne \textbf{Obligatoire} pour basculer l'\'etat.
\end{info}

% -------- Onglet Prix --------
\section{Onglet 3 --- Prix par Variable et Classe}

C'est ici que vous d\'efinissez le montant exact que chaque classe doit payer pour chaque variable. C'est le c\oe ur de la configuration financi\`ere.

\begin{etape}[title=Proc\'edure d\'etaill\'ee : D\'efinir les prix]
\begin{enumerate}[leftmargin=1.5cm,label=\textbf{\arabic*.},itemsep=12pt]
  \item \textbf{S\'electionnez l'ann\'ee} dans le menu \champ{Ann\'ee}.
  \item \textbf{S\'electionnez la classe} dans le menu \champ{Classe}. Le tableau des variables se charge.
  \item \textbf{Pour chaque variable}, saisissez le prix dans le champ num\'erique de la colonne \champ{Prix (Fbu)}.
  \item \textbf{Cliquez sur} \bouton{Enregistrer} sur la ligne correspondante.
  \item Le badge passe de \textcolor{amber}{\textit{Non d\'efini}} \`a \textcolor{emerald}{\textit{D\'efini}}.
\end{enumerate}
\end{etape}

\begin{center}
\begin{tikzpicture}
  \fill[white, rounded corners=6pt, drop shadow={shadow xshift=1pt, shadow yshift=-1pt, opacity=0.1}] (0,0) rectangle (13,3.2);
  \draw[indigo!20, rounded corners=6pt, line width=1pt] (0,0) rectangle (13,3.2);
  % Header
  \fill[slate!8, rounded corners=4pt] (0.2,2.4) rectangle (12.8,2.9);
  \node[font=\tiny\bfseries\color{slate}] at (1,2.65) {N\textdegree};
  \node[font=\tiny\bfseries\color{slate}] at (3.5,2.65) {Variable};
  \node[font=\tiny\bfseries\color{slate}] at (6.5,2.65) {Cat\'egorie};
  \node[font=\tiny\bfseries\color{slate}] at (9.5,2.65) {Prix (Fbu)};
  \node[font=\tiny\bfseries\color{slate}] at (11.8,2.65) {Action};
  % Row 1
  \node[font=\tiny\color{slate}] at (1,2) {1};
  \node[font=\tiny\bfseries\color{slate}] at (3.5,2) {Frais scolaire T1};
  \node[fill=emerald!12, rounded corners=2pt, inner sep=2pt, font=\tiny\color{emerald}] at (6,2) {\faCheck\ D\'efini};
  \node[font=\tiny\color{slate}] at (7.2,2) {Minerval};
  \node[draw=slate!20, rounded corners=3pt, minimum width=1.5cm, font=\tiny\color{slate}, inner sep=2pt] at (9.5,2) {150 000};
  \node[fill=emerald!80, rounded corners=3pt, inner sep=2pt, font=\tiny\bfseries\color{white}] at (11.8,2) {\faSave\ Enregistrer};
  % Row 2
  \node[font=\tiny\color{slate}] at (1,1.3) {2};
  \node[font=\tiny\bfseries\color{slate}] at (3.5,1.3) {Uniforme};
  \node[fill=amber!12, rounded corners=2pt, inner sep=2pt, font=\tiny\color{amber}] at (6,1.3) {\faExclamationCircle\ Non d\'efini};
  \node[font=\tiny\color{slate}] at (7.2,1.3) {Fournitures};
  \node[draw=slate!20, rounded corners=3pt, minimum width=1.5cm, font=\tiny\color{slatelight}, inner sep=2pt] at (9.5,1.3) {0};
  \node[fill=emerald!80, rounded corners=3pt, inner sep=2pt, font=\tiny\bfseries\color{white}] at (11.8,1.3) {\faSave\ Enregistrer};
  % Summary
  \node[font=\tiny\color{slate}, anchor=west] at (0.4,0.5) {\textcolor{emerald}{\faCheckCircle}\ \textbf{1} prix d\'efini \quad \textcolor{amber}{\faExclamationCircle}\ \textbf{1} en attente};
\end{tikzpicture}
\captionof{figure}{Exemple du tableau de configuration des prix}
\end{center}

\begin{attention}
Les prix sont sp\'ecifiques \`a chaque couple \textbf{ann\'ee + classe}. Si vous changez d'ann\'ee ou de classe, vous devez red\'efinir les prix. Cela permet une flexibilit\'e maximale (prix diff\'erents selon les niveaux).
\end{attention}

% -------- Onglet P\'enalit\'es --------
\section{Onglet 4 --- R\`egles de P\'enalit\'e}

Les p\'enalit\'es permettent de configurer des frais suppl\'ementaires automatiques en cas de retard de paiement. Deux types sont disponibles~:

\begin{center}
\begin{tabularx}{\textwidth}{lX}
\toprule
\textbf{Type} & \textbf{Description} \\
\midrule
Forfait & Un montant fixe ajout\'e au solde de l'\'el\`eve (ex.\ : 5\,000 Fbu de p\'enalit\'e). \\
Pourcentage & Un pourcentage du montant restant d\^u (ex.\ : 10\% du reste \`a payer). \\
\bottomrule
\end{tabularx}
\end{center}

\begin{etape}[title=Proc\'edure : Cr\'eer une r\`egle de p\'enalit\'e]
\begin{enumerate}[leftmargin=1.5cm,label=\textbf{\arabic*.},itemsep=10pt]
  \item \textbf{Cliquez sur l'onglet} \onglet{P\'enalit\'es}.
  \item \textbf{S\'electionnez l'ann\'ee} dans \champ{Ann\'ee}.
  \item \textbf{S\'electionnez la variable} concern\'ee dans \champ{Variable} (ou laissez sur \textit{Toutes} pour une p\'enalit\'e globale).
  \item \textbf{Choisissez le type}~: \textit{Forfait} ou \textit{Pourcentage (\%)}.
  \item \textbf{Saisissez la valeur} (montant en Fbu pour un forfait, pourcentage pour un \%).
  \item \textbf{Saisissez le plafond} (optionnel) --- montant maximal que la p\'enalit\'e ne peut pas d\'epasser.
  \item \textbf{Cliquez sur} \bouton{Ajouter}. La r\`egle appara\^it dans le tableau ci-dessous.
\end{enumerate}
\end{etape}

% -------- Onglet Dates Butoires --------
\section{Onglet 5 --- Dates Butoires}

Les dates butoires d\'efinissent les dates limites de paiement par variable et par classe.

\begin{etape}[title=Proc\'edure : D\'efinir une date butoire]
\begin{enumerate}[leftmargin=1.5cm,label=\textbf{\arabic*.},itemsep=10pt]
  \item \textbf{Cliquez sur l'onglet} \onglet{Dates Butoires}.
  \item \textbf{S\'electionnez l'ann\'ee} et la \textbf{classe}.
  \item \textbf{S\'electionnez la variable} concern\'ee.
  \item \textbf{Saisissez la date limite} dans \champ{Date limite}.
  \item \textbf{Cliquez sur} \bouton{Enregistrer}.
\end{enumerate}
\end{etape}

\begin{info}
Les dates d\'epass\'ees apparaissent en \textcolor{rose}{\textbf{rouge}} dans le tableau, et les dates futures en \textcolor{emerald}{\textbf{vert}}. Cela permet de rep\'erer rapidement les \'ech\'eances d\'epass\'ees.
\end{info}

% -------- Onglet Banques --------
\section{Onglet 6 --- Banques et Comptes}

Cet onglet est divis\'e en deux panneaux c\^ote \`a c\^ote~:

\begin{enumerate}[leftmargin=1.2cm]
  \item \textbf{Panneau gauche} --- \textbf{Banques}~: institutions bancaires (ex.\ : BANCOBU, KCB).
  \item \textbf{Panneau droit} --- \textbf{Comptes bancaires}~: num\'eros de compte li\'es \`a une banque.
\end{enumerate}

\begin{etape}[title=Proc\'edure : Ajouter une banque et un compte]
\begin{enumerate}[leftmargin=1.5cm,label=\textbf{\arabic*.},itemsep=10pt]
  \item \textbf{Cliquez sur l'onglet} \onglet{Banques \& Comptes}.
  \item \textbf{Dans le panneau Banques}~: saisissez le \champ{Nom} et le \champ{Sigle}, puis cliquez sur \bouton{+}.
  \item \textbf{Dans le panneau Comptes}~: s\'electionnez la \champ{Banque} dans le menu d\'eroulant, saisissez le \champ{N\textdegree\ compte}, puis cliquez sur \bouton{+}.
\end{enumerate}
\end{etape}

\begin{resultat}
Les comptes bancaires sont imm\'ediatement disponibles dans le menu d\'eroulant \champ{Compte} du formulaire de paiement (section Paiements).
\end{resultat}

\newpage
% ===================================================================
\chapter{Section 4 --- Entr\'ees et Sorties de Caisse}
% ===================================================================

\section{Objectif de cette section}

La section \textbf{Caisse} permet de suivre les flux financiers \textbf{hors paiements scolaires}~: achats de fournitures, salaires, dons re\c{c}us, frais d'entretien, etc. Elle est compl\`etement ind\'ependante du cycle de paiement des \'el\`eves.

\section{Les indicateurs de caisse}

Quatre cartes statistiques en haut~:

\begin{center}
\begin{tabularx}{\textwidth}{llX}
\toprule
\textbf{Indicateur} & \textbf{Couleur} & \textbf{Description} \\
\midrule
Total Op\'erations & Violet & Nombre total d'op\'erations enregistr\'ees \\
Total Entr\'ees & Vert & Somme de toutes les entr\'ees (recettes) \\
Total Sorties & Rose & Somme de toutes les sorties (d\'epenses) \\
Solde & Bleu & Diff\'erence Entr\'ees $-$ Sorties \\
\bottomrule
\end{tabularx}
\end{center}

\section{Pour cr\'eer une cat\'egorie d'op\'eration}

\begin{etape}[title=Proc\'edure : Cr\'eer une cat\'egorie d'op\'eration]
\begin{enumerate}[leftmargin=1.5cm,label=\textbf{\arabic*.},itemsep=10pt]
  \item \textbf{Cliquez sur l'onglet} \onglet{Cat\'egories d'op\'erations} dans la section Caisse.
  \item \textbf{S\'electionnez le type}~: \textit{Entr\'ee} ou \textit{Sortie}.
  \item \textbf{Saisissez le nom} (ex.\ : \textit{Fournitures bureau}).
  \item \textbf{Ajoutez une description} (optionnel).
  \item \textbf{Cliquez sur} \bouton{Ajouter}.
\end{enumerate}
\end{etape}

\section{Pour enregistrer une op\'eration}

\begin{etape}[title=Proc\'edure d\'etaill\'ee : Enregistrer une op\'eration de caisse]
\begin{enumerate}[leftmargin=1.5cm,label=\textbf{\arabic*.},itemsep=12pt]
  \item \textbf{Acc\'edez \`a l'onglet} \onglet{Journal des op\'erations}.
  \item \textbf{Cliquez sur} \bouton{Nouvelle op\'eration}. Le formulaire s'affiche.
  \item \textbf{S\'electionnez le type}~: \textit{Entr\'ee} ou \textit{Sortie}. Les cat\'egories correspondantes se chargent.
  \item \textbf{S\'electionnez la cat\'egorie} d'op\'eration.
  \item \textbf{Saisissez le montant} (positif obligatoirement).
  \item \textbf{Saisissez la date} de l'op\'eration.
  \item \textbf{S\'electionnez le mode} de paiement~: Esp\`eces, Ch\`eque, Virement, Autre.
  \item \textbf{Compl\'etez} les champs optionnels~: Source/B\'en\'eficiaire, R\'ef\'erence, Description.
  \item \textbf{Cliquez sur} \bouton{Enregistrer}.
\end{enumerate}
\end{etape}

\section{Pour exporter le journal}

\begin{etape}[title=Proc\'edure : Exporter les op\'erations]
\begin{enumerate}[leftmargin=1.5cm,label=\textbf{\arabic*.},itemsep=8pt]
  \item \textbf{Filtrez} les op\'erations par ann\'ee, cat\'egorie ou type si souhait\'e.
  \item \textbf{Pour un PDF}~: cliquez sur \btnrose{PDF} en haut \`a droite du journal.
  \item \textbf{Pour un Excel}~: cliquez sur \bouton{Excel} en haut \`a droite. Le fichier se t\'el\'echarge directement.
\end{enumerate}
\end{etape}

\newpage
% ===================================================================
\chapter{R\'ef\'erence rapide}
% ===================================================================

\section{Navigation entre les sections}

\begin{tabularx}{\textwidth}{llX}
\toprule
\textbf{Ic\^one Sidebar} & \textbf{Section} & \textbf{Description} \\
\midrule
\faChartPie\ (violet) & Dashboard & Tableau de bord financier avec indicateurs et graphiques \\
\faMoneyBillWave\ (vert) & Paiements & Enregistrement, validation et suivi des paiements \\
\faCogs\ (ambre) & Configuration & Cat\'egories, variables, prix, p\'enalit\'es, banques \\
\faCashRegister\ (rouge) & Caisse & Entr\'ees et sorties financi\`eres hors scolaire \\
\bottomrule
\end{tabularx}

\section{Endpoints API du module Recouvrement}

{\small
\begin{tabularx}{\textwidth}{lX}
\toprule
\textbf{Endpoint} & \textbf{Description} \\
\midrule
\texttt{/api/recouvrement/dashboard-data/} & GET --- Statistiques g\'en\'erales du dashboard \\
\texttt{/api/recouvrement/classes/} & GET --- Classes actives pour une ann\'ee \\
\texttt{/api/recouvrement/eleves/} & GET --- \'El\`eves inscrits dans une classe \\
\texttt{/api/recouvrement/variables-restant/} & GET --- Variables avec reste \`a payer d'un \'el\`eve \\
\texttt{/api/recouvrement/save-paiement/} & POST --- Enregistrer un paiement \\
\texttt{/api/recouvrement/paiements-submitted/} & GET --- Paiements en attente \\
\texttt{/api/recouvrement/paiements-validated/} & GET --- Paiements valid\'es et rejet\'es \\
\texttt{/api/recouvrement/save-categorie-variable/} & POST --- Cr\'eer une cat\'egorie \\
\texttt{/api/recouvrement/save-variable/} & POST --- Cr\'eer une variable \\
\texttt{/api/recouvrement/save-variable-prix/} & POST --- D\'efinir un prix \\
\texttt{/api/recouvrement/save-penalite/} & POST --- Cr\'eer une r\`egle de p\'enalit\'e \\
\texttt{/api/recouvrement/save-date-butoire/} & POST --- D\'efinir une date butoire \\
\texttt{/api/recouvrement/save-banque/} & POST --- Cr\'eer une banque \\
\texttt{/api/recouvrement/save-compte/} & POST --- Cr\'eer un compte bancaire \\
\texttt{/api/recouvrement/operations/} & GET --- Journal des op\'erations de caisse \\
\texttt{/api/recouvrement/save-operation/} & POST --- Enregistrer une op\'eration \\
\texttt{/api/recouvrement/invoice/<id>/} & GET --- G\'en\'erer une facture PDF \\
\bottomrule
\end{tabularx}
}

\section{Glossaire}

\begin{description}[leftmargin=1.5cm, style=nextline, font=\bfseries\color{indigo}]
  \item[Bordereau] Pi\`ece justificative (photo ou PDF) jointe \`a un paiement.
  \item[Cat\'egorie] Regroupement logique de variables de paiement.
  \item[Date butoire] Date limite de paiement pour une variable donn\'ee.
  \item[D\'erogation] Prolongation individuelle d'une date butoire pour un \'el\`eve.
  \item[Fiche classe] Document PDF r\'ecapitulatif des paiements de toute une classe.
  \item[Fiche \'el\`eve] Document PDF individuel d\'etaillant les paiements d'un \'el\`eve.
  \item[Op\'eration de caisse] Mouvement financier (entr\'ee ou sortie) hors scolaire.
  \item[P\'enalit\'e] Frais suppl\'ementaire appliqu\'e en cas de retard de paiement.
  \item[R\'eduction] Remise accord\'ee \`a un \'el\`eve (en pourcentage du prix total).
  \item[Reste \`a payer] Montant total d\^u moins les paiements valid\'es.
  \item[Variable] \'El\'ement de facturation sp\'ecifique (frais scolaire, uniforme, etc.).
\end{description}

\vspace{2cm}
\begin{center}
\begin{tikzpicture}
  \fill[coverbg, rounded corners=12pt] (0,0) rectangle (14,2.5);
  \node[text=white, font=\large\bfseries] at (7,1.5) {\faCheck\hspace{8pt}Fin du Manuel --- Module Recouvrement};
  \node[text=white!50, font=\small] at (7,0.7) {Plateforme MonEcole --- Nexora Tech --- Mai 2026 --- v1.0};
\end{tikzpicture}
\end{center}
